\section{Introduction}

Trustworthiness evaluation frameworks for large language models (LLMs) implicitly assume that safety, fairness, truthfulness, adversarial robustness, privacy, and machine ethics are orthogonal axes --- each measuring a distinct property that demands independent assessment. This assumption is operationalized in every major benchmark: TrustLLM \cite{sun2024trustllm} evaluates 16 models across 6 dimensions, HELM \cite{liang2023helm} tracks 7 metrics across 30 models, and MultiTrust \cite{zhang2024multitrust} presents 5 dimensions as a radar chart. Yet no study has asked the empirical question lurking underneath: \emph{are these dimensions actually independent?}

We show they are not. After controlling for model scale (log parameter count) and RLHF alignment status, 8 of 15 trustworthiness dimension pairs are statistically correlated at Bonferroni-corrected significance, with partial Spearman coefficients reaching $\rho = 0.971$ for the safety--privacy pair. The correlation structure is not noise --- it is a systematic geometry driven by RLHF preference training that co-optimizes 5 of 6 dimensions simultaneously, while adversarial robustness stands apart as the sole RLHF-insensitive axis.

This finding has two immediate practical implications. First, \textbf{evaluation compression}: a 3-dimension minimum evaluation set \{truthfulness, fairness, privacy\}, derived from the minimum spanning tree (MST) of the partial correlation matrix, spans the correlated cluster with 91.7\% mean bootstrap edge stability --- a 50\% reduction from the standard 6-dimension protocol. Second, \textbf{differential attention}: adversarial robustness cannot be inferred from the 5-dimension cluster and requires separate evaluation, as it follows a trajectory governed by model scale rather than alignment.

\textbf{The key insight} driving these findings is that RLHF human preference annotation creates a pervasive shared optimization signal: annotators penalize harmful, unfair, and privacy-violating outputs through the same rating judgments, simultaneously lifting safety, ethics, fairness, privacy, and truthfulness. Adversarial robustness --- which is never directly tested in preference annotation --- escapes this shared pressure and follows an independent, scale-driven path.

This reframes trustworthiness evaluation as a geometry problem rather than a dimension-counting problem. Rather than asking ``how does model X score on each of 6 dimensions?'', the correct question is ``what is the minimal set of dimensions that spans the trustworthiness correlation space?'' --- a question that existing evaluation frameworks do not address but that the MST formalism answers directly.

We make the following contributions:

\begin{enumerate}
  \item \textbf{The first, to our knowledge, systematic partial Spearman correlation analysis of LLM trustworthiness} --- We compute the full $6 \times 6$ partial Spearman correlation matrix for TrustLLM's 16-model $\times$ 6-dimension evaluation data, controlling for model scale and RLHF status, revealing a strong 2-cluster structure (silhouette = 0.637, robust across all linkage methods) not previously reported.

  \item \textbf{Robustness isolation principle} --- Adversarial robustness is empirically identified as the sole RLHF-insensitive trustworthiness dimension, with 3/3 LLaMA-2 within-family pairs showing negative robustness deltas after RLHF fine-tuning, consistent with a safety-robustness tension that is directional but attenuated at $n=16$.

  \item \textbf{RLHF joint optimization of safety and ethics} --- A within-family natural experiment on LLaMA-2 (7B, 13B, 70B) demonstrates that RLHF fine-tuning simultaneously increases safety ($\Delta = +0.63$ average) and machine ethics ($\Delta = +0.42$ average) at every scale, providing mechanistic evidence for the 5-dimension RLHF-shaped cluster.

  \item \textbf{MST-derived minimum evaluation set} --- The minimum spanning tree identifies \{truthfulness, fairness, privacy\} as a sufficient 3-dimension evaluation set with mean bootstrap edge frequency 0.917 \cite{tumminello2005tool}, enabling principled evaluation compression.
\end{enumerate}

We organize the paper as follows: Section~\ref{sec:related} surveys related work on trustworthiness benchmarking and benchmark correlation analysis. Section~\ref{sec:method} describes our partial correlation and clustering methodology. Section~\ref{sec:setup} details the experimental setup. Section~\ref{sec:results} presents results. Section~\ref{sec:discussion} discusses implications and limitations. Section~\ref{sec:conclusion} concludes.
